% Interior house style: a readable serif text with a quiet sans-serif hierarchy.
% Keep linguistic semantics and content macros in grammar.tex.
\newif\ifbookmainmatter
\AddToHook{cmd/mainmatter/after}{\bookmainmattertrue}
\AddToHook{cmd/backmatter/after}{\bookmainmatterfalse}
\renewcommand{\thepart}{\Roman{part}}
\renewcommand{\thesection}{\thechapter.\arabic{section}}
\renewcommand{\thetable}{\thechapter.\arabic{table}}
\renewcommand{\thefigure}{\thechapter.\arabic{figure}}

\newcommand{\booktitlepage}{%
  \begin{titlepage}
    \thispagestyle{empty}
    \raggedright
    {\sffamily\large Roel Christian Yambao\par}
    \vspace*{35mm}
    {\sffamily\bfseries\fontsize{32}{38}\selectfont
      A Reference Grammar\par of the Iridian Language\par}
    \vspace{10mm}
    {\Large\itshape Ircevni kostgramatik\par}
    \vfill
    {\sffamily\small First edition\quad 2026\par}
  \end{titlepage}%
}

\makeatletter
\def\@makechapterhead#1{%
  \thispagestyle{plain}%
  \vspace*{12mm}%
  {\parindent\z@\raggedright\normalfont
    {\sffamily\small\MakeUppercase{\chaptername}\enspace\thechapter\par}
    \vspace{5mm}%
    {\sffamily\bfseries\fontsize{27}{33}\selectfont #1\par}
    \vspace{11mm}}}
\def\@makeschapterhead#1{%
  \thispagestyle{plain}%
  \vspace*{12mm}%
  {\parindent\z@\raggedright\normalfont
    {\sffamily\bfseries\fontsize{27}{33}\selectfont #1\par}
    \vspace{11mm}}}

\renewcommand\section{\@startsection{section}{1}{\z@}%
  {-1.7\baselineskip plus -.3\baselineskip minus -.2\baselineskip}%
  {.65\baselineskip}%
  {\normalfont\sffamily\bfseries\large\raggedright}}
\renewcommand\subsection{\@startsection{subsection}{2}{\z@}%
  {-1.25\baselineskip plus -.2\baselineskip minus -.1\baselineskip}%
  {.45\baselineskip}%
  {\normalfont\bfseries\normalsize\raggedright}}
\renewcommand\subsubsection{\@startsection{subsubsection}{3}{\z@}%
  {-1\baselineskip plus -.2\baselineskip minus -.1\baselineskip}%
  {.35\baselineskip}%
  {\normalfont\itshape\normalsize\raggedright}}

% Reserve enough label width for identifiers such as 3.12 and 10.1.
\def\l@table{\@tocline{0}{3pt plus2pt}{0pt}{3em}{}}
\let\l@figure\l@table

% Compact chapter identifiers and aligned section numbers in the contents.
\renewcommand{\tocchapter}[3]{\indentlabel{\@ifnotempty{#2}{#2\quad}}#3}
\renewcommand{\chaptermark}[1]{\markboth{\ifbookmainmatter\thechapter\quad\fi #1}{#1}}
\renewcommand{\sectionmark}[1]{\markright{\thesection\quad #1}}
\makeatother

\fancypagestyle{plain}{%
  \fancyhf{}
  \fancyfoot[LE,RO]{\sffamily\small\thepage}
  \renewcommand{\headrulewidth}{0pt}
  \renewcommand{\footrulewidth}{0pt}}
\pagestyle{fancy}
\fancyhf{}
\fancyhead[LE]{\parbox[b]{\textwidth}{\sffamily\footnotesize\raggedright\nouppercase{\leftmark}}}
\fancyhead[RO]{\parbox[b]{\textwidth}{\sffamily\footnotesize\raggedleft\nouppercase{\rightmark}}}
\fancyfoot[LE,RO]{\sffamily\small\thepage}
\renewcommand{\headrulewidth}{0pt}
\renewcommand{\footrulewidth}{0pt}
\setcounter{secnumdepth}{2}
\setcounter{tocdepth}{1}
\renewcommand{\contentsname}{Contents}
